% Generated draft: requires mathematical and source review.
\documentclass[11pt]{article}
\usepackage[margin=25mm]{geometry}
\usepackage{fontspec}
\setmainfont{Latin Modern Roman}
\usepackage{amsmath,amssymb,mathtools,unicode-math}
\setmathfont{Latin Modern Math}
\usepackage{enumitem,needspace,xurl,hyperref}
\setlength{\parindent}{0pt}
\setlength{\parskip}{5pt}
\setlength{\emergencystretch}{3em}
\newcommand{\N}{\mathbb{N}}
\newcommand{\Z}{\mathbb{Z}}
\newcommand{\Q}{\mathbb{Q}}
\newcommand{\R}{\mathbb{R}}
\newcommand{\C}{\mathbb{C}}
\newcommand{\F}{\mathbb{F}}
\newcommand{\A}{\mathbb{A}}
\newcommand{\PP}{\mathbb{P}}
\newcommand{\E}{\mathbb{E}}
\newcommand{\eps}{\varepsilon}
\newcommand{\dd}{\,\mathrm{d}}
\begin{document}
\section*{UnsolvedMath ID: 800010}
\textbf{Source problem code:} catalogue entry
\subsection{Definitions and mathematical statement}
The supplied problem, attributed to E. S. Lyapin, asks: under what conditions is a system of identical inequalities (also called identical relations of precedence in the source) equivalent to a system of identities?

The source does not define these inequalities, their underlying structures, or the intended meaning of equivalence. The following schema makes the missing parameters explicit without fixing them. Let \(L\) be a signature with a designated binary relation symbol \(\preceq\), and let \(\mathcal K\) be a specified class of \(L\)-structures representing the intended semigroup setting. Both \(L\) and \(\mathcal K\), including all axioms for \(\preceq\), remain unspecified by the supplied record. If an identical inequality means a universally quantified relation between terms, write a system as
\[
\Sigma=\left\{\forall\mathbf x_a\;\bigl(u_a(\mathbf x_a)\preceq v_a(\mathbf x_a)\bigr):a\in A\right\},
\]
where each \(\mathbf x_a\) contains all variables occurring in the displayed terms. A system of identities has the form
\[
T=\left\{\forall\mathbf y_b\;\bigl(p_b(\mathbf y_b)=q_b(\mathbf y_b)\bigr):b\in B\right\}.
\]
Under the interpretation that equivalence means satisfaction by exactly the same structures in \(\mathcal K\), the requested characterization would concern those \(\Sigma\) for which there exists such a \(T\) satisfying
\[
\forall M\in\mathcal K,\qquad M\models\Sigma\ \Longleftrightarrow\ M\models T.
\]
This is a conditional formulation, not a determination of the source's intended semantics. The permitted signatures, background class, meaning of precedence, restrictions on the index sets \(A,B\), and whether equivalence instead concerns reducts or another construction require clarification.

The supplied source states that the case of a system consisting of a single identical inequality is solved in the reference labeled \([\mathrm{Mak1}]\). Its hypotheses and conclusion are not supplied, so this statement does not specify a theorem that can be incorporated here.
\subsection{Short English statement}
Lyapin asks when a system of universally interpreted inequalities or precedence relations can be replaced by identities with equivalent meaning. The supplied record does not specify enough semantics to determine a unique precise formulation. It labels the general problem open and reports the single-inequality case solved in \([\mathrm{Mak1}]\); these are supplied source assessments, not independently verified conclusions.
\subsection{Sources}
\begin{itemize}
\item Local UnsolvedMath registry, permanent ID \(800010\), exact source problem code \texttt{catalogue entry}, title \emph{The Lyapin's Notebook --- Problem 10}. The supplied registry status is \texttt{open}.
\item Bershadsky, Kublanovsky, and Mashevitzky, \emph{The Lyapin's Notebook: Unsolved Problems in Semigroup Theory} (2026), Problem 10, attributed within the statement to E. S. Lyapin. Supplied source URL: \url{https://arxiv.org/abs/2604.04763}. The record reports access on 2026-10-16 and extraction by \texttt{arxiv\_latex}; these provenance details are supplied and have not been independently verified.
\item The source cites \([\mathrm{Mak1}]\) for the assertion that systems consisting of a single identical inequality have been treated. No bibliographic details or link for this reference are supplied.
\end{itemize}
\end{document}
